\documentclass[11pt,a4paper]{article}
\usepackage[margin=25mm]{geometry}
\usepackage{amsmath,amssymb,amsthm,hyperref}
\setlength{\emergencystretch}{2em}
\newtheorem{lemma}{Lemma}
\newtheorem{theorem}[lemma]{Theorem}
\newcommand{\gen}[1]{\langle#1\rangle}
\title{Faithful symmetric lifts on pairs}
\author{Human Compression project: a component of the L044 recovery}
\date{}
\begin{document}\maketitle
Let $G=A_{2m}=\gen{\rho_0,\ldots,\rho_{r-1}}$ be naturally
permuted and have a string C-marking. Suppose its maximal
omission $H=G_i$ is transitive on $m\ge8$ pairs, with block
image $S_m$. Write $V=\mathbb F_2^m$, let $z$ be the constant
flip, and let $E$ be the even-weight subspace. Choose the
minimal distinguished block-generating family $L$, with
$B=\gen L$, and the commuting and remaining families $C,R$.
We acquire the full L047 decomposition and the accepted actual
pair-centralizer lemma. In any putative case $r\ge m$, $L$
is an interval, $|C|\le1$ and $|R|\le1$.

\begin{lemma}\label{onefixed}
Suppose a permutation group preserves one pair setwise and
induces $S_t$ on its other $t\ge4$ pairs. Every permutation
centralizing that group preserves all its pairs.
\end{lemma}
\begin{proof}
Its point orbits on the other pairs have size at least $t$,
whereas those in the distinguished pair have size at most two.
A centralizer preserves these two unions of orbits.
On the $2t$ points the acquired pair-centralizer lemma applies:
the block group $S_t$ is primitive and is not a regular
prime-dihedral group of order $2t$ for $t\ge4$. It forces
preservation of every pair there. The remaining two points
are the distinguished pair. No faithfulness of the original
group's restriction to either union is presumed; the lemma
is applied to its actual induced permutation group.
\end{proof}

\begin{lemma}\label{core}
A connected string C-marking of the natural $S_m$, of rank
$m-2$, has, up to simultaneous conjugacy and reversal, the form
\[
 (2\ 3),\quad(1\ 2)(3\ 4),\quad(4\ 5),\quad
 (5\ 6),\ldots,(m-1\ m),\qquad m\ge8.                 \tag{1}
\]
Its tail-end omission is the natural $S_{m-1}$ fixing $m$;
its core-end omission is $S_{\{1,2\}}\times S_{\{3,\ldots,m\}}$.
\end{lemma}
\begin{proof}
For $m\ge9$ we acquire the complete B5 classification,
including its general four-point cores and attached trees.
Its core B produces a triangle of noncommutation with the
first attached edge, and therefore has no string order.
In core C, that first attached edge fails to commute with
both core marks and with the next attached edge; its degree
in the noncommutation graph is at least three. There is
such a next edge because the root is a leaf and the tail
has at least two edges. This also cannot be a string.
In core A, the first two marks are the first two entries
of (1). The attached tree has no vertex of degree three,
since its incident transpositions would give a triangle
of noncommutation. Hence it is a path. The resulting
noncommutation diagram is itself a path, so its only string
orders are its path order and reversal, giving (1).

The degree-eight boundary is a complete new execution of
the acquired L047 prefix-stabilizer engine, with the actual
ambient group changed to $S_8$ and target and depth both six.
The search includes every nonidentity involution, requires
independence and the string commutations, and requires
noncommutation at each consecutive position. It takes
orbits only under the stabilizer of the entire prefix.
Every discarded prefix either already has a redundant mark,
fails a necessary end-intersection equality, or has
insufficient subgroup index or centralizer quotient index
for the remaining independent marks. These are exactly
the retained completeness conditions of that engine.
In particular, the centralizer-quotient test is applied
only to a possible completion with full IP: the subgroup
on the future marks, except the immediate next mark,
intersects the old prefix group trivially and so injects
into its centralizer modulo the old center.

All rank-six candidates are tested against the full
intersection property. The complete run has prefix counts
\[
                      1,4,63,115,22,9,3
\]
and exactly two surviving ordered conjugacy classes,
with types $[4,6,3,3,3]$ and $[3,3,3,6,4]$.
A separate exact check verifies generation of literal $S_8$,
all involutions and string commutations, all $64^2$ subset
intersections for each tuple, and simultaneous conjugacy
to (1) or its reverse. It also checks that these two
ordered classes differ. Thus the same conclusion holds
at eight. This is only the stated rank-six connected
boundary, not a census of every rank or every C-group.

The tail omission in (1) fixes $m$ and generates $S_{m-1}$:
the core with its first attachment generates $S_5$, and
each further transposition attaches one new point.
For the other omission, if $u=(1\ 2)(3\ 4)$ and
$v=(4\ 5)$, then $(uv)^3=(1\ 2)$. This recovers
$(3\ 4)$ and all of the remaining path transpositions.
The generated group is exactly the displayed product.
\end{proof}

\begin{lemma}\label{lift}
A subgroup of $C_2\wr S_m$ projecting isomorphically to
$S_m$, for $m\ge8$, is, after coordinates are chosen in
the pairs, either the ordinary doubled action or
\[
       \pi:(j,\epsilon)\longmapsto
                   (j^\pi,\epsilon+\operatorname{sgn}_2\pi).
                                                               \tag{2}
\]
The second action is excluded for a subgroup of $A_{2m}$
when $m$ is odd.
\end{lemma}
\begin{proof}
If the subgroup has two point orbits, each meets each pair
once; these two rows give the ordinary doubled action.
Otherwise it is transitive, and a point stabilizer has
index two in a block stabilizer $S_{m-1}$. Its unique
index-two subgroup is $A_{m-1}$. The coset action on
$S_m/A_{m-1}$ is exactly (2); an equivariant identification
can be chosen to match the specified pairs.
In (2) an odd $\pi$ is its ordinary doubled permutation,
which is even, multiplied by the flip of all $m$ pairs.
For odd $m$ this is an odd permutation, as claimed.
\end{proof}

\begin{theorem}\label{main}
Under the standing hypotheses, a putative rank $r\ge m$
can have neither $|L|=m-1$ nor $B\cap V=1$.
Thus it forces
\[
 r=m,\quad |L|=m-2,\quad C=\varnothing,\quad |R|=1,
 \quad B\cap V=\gen z,\quad H\cap V=E,
 \quad m\text{ even}.                                \tag{3}
\]
\end{theorem}
\begin{proof}
The invariant subspaces of the permutation module are
$0,\gen z,E,V$, and $H\cap V\le E$. If $C$ is nonempty,
its mark is $z$ and $m$ is even. IP excludes $z$ from $B$,
so the same subspace list gives $B\cap V=1$.
If $R$ were also nonempty, $\gen{L,R}$ would strictly
contain $B$ and thus have nonzero kernel. For even $m$
that kernel contains $z$, contrary to IP with the mark
in $C$. Consequently
\[
 |C|+|R|\le1,\qquad r\le |L|+2,\qquad
                         m-2\le |L|\le m-1.          \tag{4}
\]

If $|L|=m-1$, its independent projected involutory
string is the Coxeter path, by the complete acquired
extremal theorem. The sign-modified alternative in that
theorem cannot consist of involutions: an edge adjacent
to its distinguished tree edge has order three.
The omitted $\rho_i$ commutes with the lifted subgroup
obtained by removing the adjacent endpoint of $L$.
That subgroup fixes one original pair setwise and induces
$S_{m-1}$ on the rest. Lemma~\ref{onefixed} makes
$\rho_i$ preserve all pairs, contradicting the parent
$A_{2m}$.

It remains that $|L|=m-2$, $r=m$, and exactly one mark
lies outside $L$ in $H$. Suppose first that $B\cap V=1$.
Its faithful projection is a connected C-marking of $S_m$:
$L$ is an interval of the connected parent diagram, and
faithfulness preserves its consecutive noncommutations.
Lemma~\ref{core} applies.
If the omitted $\rho_i$ is next to the tail end, it
centralizes the lifted tail omission and is excluded by
Lemma~\ref{onefixed}. We may therefore orient (1) with
$\rho_i$ next to its core end. It centralizes the lifted
core omission $S_2\times S_b$, where $b=m-2$.
The pure $A_b$ on the larger block orbit fixes exactly
the four points over the small orbit in either action
of Lemma~\ref{lift}. Thus $\rho_i$ preserves those
four points and the other $2b$ points separately.

If the one extra mark is $C=\{z\}$, then $m$ and $b$
are even. In the ordinary doubled action the centralizer
on the $2b$ points is $\gen{z_b}$, whose elements are
even. On the four points the core omission acts as a
block-interchanging double transposition. Its even
centralizer is the Klein four group, preserving the
original pairs. Since $\rho_i$ is even, its two restrictions
have this parity and it preserves every pair.
In the action (2), an odd element of the large $S_b$
induces the full flip on the four small points, while
the small $S_2$ induces the other block-interchanging
double transposition. Their images generate the regular
Klein four group there, whose centralizer is itself and
preserves the original pairs. On the $2b$ points the
pair-centralizer lemma, with block image $S_b$, again
forces preservation of the pairs. This case is also
impossible.

If instead the extra mark is in $R$, it is adjacent to
the opposite tail end of $L$, because the omitted label
occupies the core neighbour. It centralizes the lifted
tail omission. Its block image centralizes the natural
$S_{m-1}$ fixing the tail point, hence is trivial.
It is a flip vector constant off that endpoint.
If $m$ is even, evenness forces this vector to be $z$,
which commutes with $L$ and would belong to $C$, not $R$.
Thus $m$ is odd, the action of $B$ is ordinary doubled,
and the extra mark is $z+e_m$, flipping every pair
except the tail pair.
On the four core points it is their full flip, whereas
the core omission supplies the block-interchanging
double transposition. These generate a regular Klein
four group. Since the omitted mark commutes with both,
its restriction preserves the two original pairs there.
On the remaining points it centralizes the doubled
$S_b$, whose centralizer is $\gen{z_b}$ and preserves
all its pairs. This again makes the parent imprimitive.
We have excluded every faithful $B$.

Finally the module list leaves $B\cap V=\gen z$ or $E$.
The latter would give $B=H$, contrary to the extra
independent mark. Thus the kernel is $\gen z$ and $m$
is even. The strict inclusion $B<H$ forces $H\cap V=E$.
IP excludes $C=\{z\}$, so the extra mark belongs to $R$.
Together with (4) this gives every assertion in (3).
\end{proof}

\paragraph{Attribution and finite premises.}
The rank route is source Section~3.2. The L047 decomposition,
complete symmetric extremal theorem and B5 core classification
remain acquired in their full reviewed scopes. The new work
provides actual-cover centralizers, the parity argument for
both faithful lifts, and the complete degree-eight string
boundary. Its GAP engine changes only the ambient group and
one target specification in the acquired L047 code. The guarded
checker rejects optimized Python, clears that variable for
children, pins inputs and checks actual tuples separately.
GAP exact operations remain software premises. The nontrivial
kernel case (3) is an explicit remaining obligation; no
whole-original, open-question, Lean or compression credit
is claimed by this component.
\end{document}
